\documentclass[a4paper, serif, 10pt]{moderncv}

%% moderncv
% style and color
\moderncvstyle{banking}
\moderncvcolor{black}

% renew moderncv command to adapt for CJK colon
\renewcommand*{\cvitem}[3][.25em]{%
  \ifstrempty{#2}{}{\hintstyle{#2}：}{#3}%
  \par\addvspace{#1}}

\renewcommand*{\cvitemwithcomment}[4][.25em]{%
  \savebox{\cvitemwithcommentmainbox}{\ifstrempty{#2}{}{\hintstyle{#2}：}#3}%
  \setlength{\cvitemwithcommentmainlength}{\widthof{\usebox{\cvitemwithcommentmainbox}}}%
  \setlength{\cvitemwithcommentcommentlength}{\maincolumnwidth-\separatorcolumnwidth-\cvitemwithcommentmainlength}%
  \begin{minipage}[t]{\cvitemwithcommentmainlength}\usebox{\cvitemwithcommentmainbox}\end{minipage}%
  \hfill% fill of \separatorcolumnwidth
  \begin{minipage}[t]{\cvitemwithcommentcommentlength}\raggedleft\small\itshape#4\end{minipage}%
  \par\addvspace{#1}}

%% page layout/margins
\usepackage[top=2.5cm, bottom=2.5cm, left=1.5cm, right=1.5cm]{geometry}

\nopagenumbers{}

%% Babel config for English language
\usepackage[english]{babel}

%% fontspec
\usepackage{fontspec}

\IfFontExistsTF{Linux Libertine}{
  \setmainfont[Ligatures={TeX, Common}, Numbers=Lining, ItalicFont=Linux Libertine]{Linux Libertine}
}{}
\IfFontExistsTF{Linux Libertine O}{
  \setmainfont[Ligatures={TeX, Common}, Numbers=Lining, ItalicFont=Linux Libertine O]{Linux Libertine O}
}{}

%% CTeX
% CJK support, used to show CJK characters in the resume
%
% - fontset=none: disable builtin fontset but instead set the CJK font manually
% - heading=false: disable ctex heading
% - punct=kaiming: use kaiming punctuations styles for CJK
% - scheme=plain: use plain scheme, do not override \normalsize font size
% - space=auto: space settings for CJK characters
%
% ref:
% - http://ctan.mirrorcatalogs.com/language/chinese/ctex/ctex.pdf
\usepackage[UTF8, heading=false, punct=kaiming, scheme=plain, space=auto]{ctex}

\IfFontExistsTF{Noto Serif CJK SC}{
  \setCJKmainfont{Noto Serif CJK SC}
}{}
\IfFontExistsTF{Noto Sans CJK SC}{
  \setCJKsansfont{Noto Sans CJK SC}
}{}

%% line spacing
\usepackage{setspace}
\setstretch{1.125}

%% URL styling - use normal text instead of monospace
\urlstyle{same}

% Auto-underline all links
% Defer until \begin{document} so hyperref is already loaded
\AtBeginDocument{%
  \let\oldhref\href
  \renewcommand{\href}[2]{\oldhref{#1}{\underline{#2}}}%
}

%% Patch \httplink and \httpslink to handle full URLs
% The original moderncv macros always prepend http:// or https:// to URLs.
% This causes issues when the URL already has a protocol (e.g., https://example.com)
% resulting in malformed URLs like https://https://example.com.
% This patch checks if the URL already contains "://" and uses it directly if so.
% Note: We use \str_set:Nx (x-expansion) to expand macros like \@homepage first.
\ExplSyntaxOn
\str_new:N \g_yamlresume_url_str

\RenewDocumentCommand{\httpslink}{O{}m}{
  \str_set:Nx \g_yamlresume_url_str {#2}
  \str_if_in:NnTF \g_yamlresume_url_str {://}
    { \url{#2} }
    { \url{https://#2} }
}

\RenewDocumentCommand{\httplink}{O{}m}{
  \str_set:Nx \g_yamlresume_url_str {#2}
  \str_if_in:NnTF \g_yamlresume_url_str {://}
    { \url{#2} }
    { \url{http://#2} }
}
\ExplSyntaxOff

\name{林曉彤}{}
\title{資深用戶體驗設計師，專注金融科技及流動應用程式}
\phone[mobile]{+852 9123 4567}
\email{hello@lamhiutung.design}
\homepage{https://www.lamhiutung.design}

\address{中國香港，香港島，香港，中環皇后大道中 100 號}{}{}

\extrainfo{{\small \faBehance}\ \href{https://www.behance.net/lamhiutung}{@lamhiutung} {} {} {} • {} {} {}
{\small \faLinkedin}\ \href{https://www.linkedin.com/in/lamhiutung}{@lam-hiu-tung} {} {} {} • {} {} {}
{\small \faDribbble}\ \href{https://dribbble.com/lamhiutung}{@lamhiutung}}

\begin{document}

\maketitle

\section{簡介}

\cvline{}{具超過六年經驗的用戶體驗設計師，擅長為金融科技及流動應用程式提供以用戶為中心的設計方案。
%
\begin{itemize}
\item 主導大型銀行應用程式的用戶體驗重塑，提升用戶滿意度達 35\%
\item 熟悉設計系統建立及管理，確保產品一致性及開發效率
\item 善於透過用戶研究及可用性測試，發掘用戶需要並轉化為具體設計方案
\item 與產品經理及工程師緊密合作，推動設計主導的產品開發流程
\end{itemize}}
\section{教育背景}

\cventry{2015年9月–2019年6月}
        {學士，設計學（互動媒體），成績：3.5}
        {香港理工大學}
        {\url{https://www.polyu.edu.hk/}}
        {}
        {\begin{itemize}
\item 專攻互動媒體設計，畢業作品獲學院頒發「最佳設計項目」獎
\item 透過多個跨學科項目，累積用戶研究、原型設計及可用性測試的實戰經驗
\item 積極參與設計學會活動，並擔任學生設計展覽的策展人
\end{itemize}
\textbf{課程}：用戶體驗設計、人機互動、視覺傳達設計、設計研究方法、流動應用程式設計、設計心理學}
\section{工作經歷}

\cventry{2022年7月至今}
        {高級用戶體驗設計師}
        {Fintastic 金融科技有限公司}
        {\url{https://www.fintastic.example}}
        {}
        {\begin{itemize}
\item 主導公司旗艦投資應用程式的設計策略，與產品及工程團隊合作推出新功能
\item 建立及維護設計系統，統一組件規範，縮短設計到開發的交付時間約 30\%
\item 規劃及執行季度用戶研究計劃，包括深度訪談及可用性測試，為產品路線圖提供建議
\item 指導初級設計師，並定期舉辦設計分享會，促進團隊專業成長
\end{itemize}
\textbf{關鍵字}：設計系統、用戶研究、可用性測試、金融科技、團隊領導}

\cventry{2019年9月–2022年6月}
        {用戶體驗設計師}
        {Digital Gear Studio}
        {\url{https://www.digitalgear.example}}
        {}
        {\begin{itemize}
\item 負責多個客戶項目的用戶體驗設計，涵蓋零售、旅遊及教育等行業
\item 運用設計思考方法，與客戶共同定義問題，並製作線框圖及互動原型
\item 進行競爭分析及用戶訪談，將研究結果轉化為清晰的用戶旅程地圖
\item 與視覺設計師及開發人員協作，確保設計在技術上可行並符合品牌形象
\end{itemize}
\textbf{關鍵字}：設計思考、用戶旅程地圖、互動原型、客戶溝通}

\cventry{2018年6月–2018年8月}
        {用戶體驗設計實習生}
        {雲創科技}
        {\url{https://www.cloudinnovation.example}}
        {}
        {\begin{itemize}
\item 協助高級設計師進行用戶訪談及問卷調查，整理研究報告
\item 參與設計工作坊，製作低保真原型及可用性測試材料
\item 與開發團隊合作，撰寫設計規格文件及交接材料
\end{itemize}
\textbf{關鍵字}：用戶訪談、問卷調查、低保真原型}
\section{語言}

\cvline{粵語}{母語或雙語水平 \hfill \textbf{關鍵字}：母語}
\cvline{英語}{完全專業水平 \hfill \textbf{關鍵字}：IELTS 8.0、商業英語}
\cvline{普通話}{完全專業水平 \hfill \textbf{關鍵字}：普通話水平測試二級甲等}
\section{專業技能}

\cvline{用戶研究}{專家 \hfill \textbf{關鍵字}：深度訪談、可用性測試、問卷調查、人物誌、用戶旅程地圖}
\cvline{原型與設計工具}{專家 \hfill \textbf{關鍵字}：Figma、Sketch、Adobe XD、InVision、Protopie}
\cvline{互動設計}{高級 \hfill \textbf{關鍵字}：線框圖、微互動、無障礙設計、設計系統}
\cvline{前端協作}{初級 \hfill \textbf{關鍵字}：HTML、CSS、JavaScript}
\section{榮譽}

\cventry{2022年10月}
        {香港設計師協會}
        {香港設計師協會環球設計大獎 2022 – 銀獎}
        {}
        {}
        {憑藉「智能銀行應用程式用戶體驗重塑」項目，獲得數碼介面組別銀獎，表揚項目在用戶體驗及視覺設計上的卓越表現。}
\section{證書}

\cventry{2021年1月}
        {Human Factors International}
        {Certified Usability Analyst}
        {\url{https://www.humanfactors.com/}}
        {}
        {}
\section{出版刊物}

\cventry{2023年3月}
        {「建立可持續設計系統的五大原則」}
        {Medium（UX Collective）}
        {\url{https://medium.com/ux-collective/design-system-principles}}
        {}
        {分享在金融科技公司建立及管理設計系統的實際經驗，包括組件設計、文件規範及團隊協作流程。}
\section{項目}

\cventry{2022年1月–2022年12月}
        {為一家主流銀行重新設計流動應用程式，整合帳戶管理、投資及支付功能，提升整體用戶體驗。
}
        {智能銀行應用程式重塑}
        {\url{https://www.fintastic.example/projects/banking-app}}
        {}
        {\begin{itemize}
\item 透過用戶研究及數據分析，辨識用戶痛點，並制定設計策略
\item 設計全新的資訊架構及互動流程，簡化複雜交易步驟
\item 建立可擴展的設計系統，支援 iOS 及 Android 平台
\item 推出後應用程式商店評分由 3.2 升至 4.5 星
\end{itemize}
\textbf{關鍵字}：金融科技、設計系統、用戶研究、可用性測試}
\section{興趣愛好}

\cvline{行山}{麥理浩徑、露營、攝影}
\cvline{烹飪}{烘焙、港式茶餐廳美食}
\section{志願服務}

\cventry{2023年1月至今}
        {網站設計義工}
        {香港導盲犬協會}
        {\url{https://www.guidedogs.org.hk/}}
        {}
        {\begin{itemize}
\item 協助協會重新設計官方網站，改善資訊架構及捐款流程
\item 拍攝及挑選照片，提升網站視覺吸引力及故事性
\end{itemize}}

\end{document}